% Options for packages loaded elsewhere
\PassOptionsToPackage{unicode}{hyperref}
\PassOptionsToPackage{hyphens}{url}
%
\documentclass[
  11pt,
]{article}
\usepackage{lmodern}
\usepackage{amssymb,amsmath}
\usepackage{ifxetex,ifluatex}
\ifnum 0\ifxetex 1\fi\ifluatex 1\fi=0 % if pdftex
  \usepackage[T1]{fontenc}
  \usepackage[utf8]{inputenc}
  \usepackage{textcomp} % provide euro and other symbols
\else % if luatex or xetex
  \usepackage{unicode-math}
  \defaultfontfeatures{Scale=MatchLowercase}
  \defaultfontfeatures[\rmfamily]{Ligatures=TeX,Scale=1}
\fi
% Use upquote if available, for straight quotes in verbatim environments
\IfFileExists{upquote.sty}{\usepackage{upquote}}{}
\IfFileExists{microtype.sty}{% use microtype if available
  \usepackage[]{microtype}
  \UseMicrotypeSet[protrusion]{basicmath} % disable protrusion for tt fonts
}{}
\makeatletter
\@ifundefined{KOMAClassName}{% if non-KOMA class
  \IfFileExists{parskip.sty}{%
    \usepackage{parskip}
  }{% else
    \setlength{\parindent}{0pt}
    \setlength{\parskip}{6pt plus 2pt minus 1pt}}
}{% if KOMA class
  \KOMAoptions{parskip=half}}
\makeatother
\usepackage{xcolor}
\IfFileExists{xurl.sty}{\usepackage{xurl}}{} % add URL line breaks if available
\IfFileExists{bookmark.sty}{\usepackage{bookmark}}{\usepackage{hyperref}}
\hypersetup{
  hidelinks,
  pdfcreator={LaTeX via pandoc}}
\urlstyle{same} % disable monospaced font for URLs
\usepackage[margin=1in]{geometry}
\setlength{\emergencystretch}{3em} % prevent overfull lines
\providecommand{\tightlist}{%
  \setlength{\itemsep}{0pt}\setlength{\parskip}{0pt}}
\setcounter{secnumdepth}{-\maxdimen} % remove section numbering
\usepackage{esint}
\setmathfont{Asana Math}
\setlength{\emergencystretch}{3em}

\author{}
\date{}

\begin{document}

\hypertarget{low-energy-matched-caps-and-gauge-free-first-order-stability}{%
\section{Low-energy matched caps and gauge-free first-order
stability}\label{low-energy-matched-caps-and-gauge-free-first-order-stability}}

28 September 2026. Continuation of I44 and I52, after the marked-center
construction. All statements below concern actual smooth metrics.

The local construction fills a nearly Euclidean annulus at a cost
controlled by its full curvature energy and its zeroth-order deviation.
It preserves the original outer collar exactly. Applied on sufficiently
long annular ranges, it removes the gauge and scalar-curvature
hypotheses from the lower first-order bound near selected recovery
metrics. It does not parametrize arbitrary near-minimizers.

\hypertarget{a-matched-cap-from-zeroth-order-closeness-and-critical-energy}{%
\subsection{1. A matched cap from zeroth-order closeness and critical
energy}\label{a-matched-cap-from-zeroth-order-closeness-and-critical-energy}}

Let \(e\) denote the Euclidean metric and
\(\mathcal A=\{1<|x|<4\}\subset\mathbb R^4\). The curvature norm can use
either standard tensor convention; the universal constants absorb that
fixed choice.

\textbf{Theorem 1 (small-energy matched cap).} There are
\(a_0,\varepsilon_0,C>0\) such that every smooth metric \(g\) on
\(\mathcal A\) with \[
 a=\|g-e\|_{L^\infty(e)}\le a_0,\qquad
 E_0=\int_{\mathcal A}|\operatorname{Rm}_g|^2\,dV_g
          \le\varepsilon_0
\] has a smooth extension \(G\) to \(B_4\), equal to \(g\) on
\(\{7/2<|x|<4\}\), satisfying \[
 \int_{B_4}|\operatorname{Rm}_G|^2\,dV_G
                         \le C(E_0+a^2).                    \tag{1.1}
\] There is no derivative bound on the input. No boundary diffeomorphism
is needed. In particular (1.1) bounds the self-dual Weyl cost.

The proof uses a localized Ricci--DeTurck deformation. The metric is not
required to solve any stationary equation. The small \(C^0\) hypothesis
supplies the parabolic and Sobolev control; we do not invoke
curvature-only harmonic-coordinate regularity.

\hypertarget{regularization-and-a-uniform-smoothing-time}{%
\subsubsection{1.1 Regularization and a uniform smoothing
time}\label{regularization-and-a-uniform-smoothing-time}}

Embed \(B_5\) in a fixed flat torus. Extend \(g\) as
\(g_0=e+\zeta(g-e)\), where \(\zeta\) is supported in \(\mathcal A\) and
equals one on \[
 \Omega=\{9/8<|x|<15/4\}.
\] This is a smooth positive extension with the same small supremum
bound. Its curvature away from \(\Omega\) can be arbitrarily large; no
estimate below charges that curvature.

Choose a fixed smooth radial cutoff with \[
 \operatorname{supp}\chi\subset\{5/4<|x|<13/4\},\qquad
 \chi=1\ \text{on }\{7/4<|x|<9/4\},\qquad 0\le\chi\le1 .
\] For \(0<\eta<1\), put \(q=\chi^2+\eta^2\), and solve \[
 \partial_tg_\eta=-2q\,\operatorname{Ric}(g_\eta)
                    +\mathcal L_{qW(g_\eta)}g_\eta,\qquad
 W^k(g)=g^{ab}\Gamma(g)^k_{ab},\qquad g_\eta(0)=g_0.   \tag{1.2}
\] The background connection is the flat one. The full Lie derivative
\(\mathcal L_{qW}\), not \(q\mathcal L_W\), is important both for the
coordinate equation and the geometric energy identity. For every fixed
\(\eta>0\), this is a smooth strictly parabolic system. Local existence
is the classical DeTurck reduction and quasilinear parabolic existence
theorem.

Here is the uniform continuation estimate needed for the proof. Write
\(u=g_\eta-e\). In flat coordinates (1.2) takes the form \[
 \partial_tu=q\,g_\eta^{ab}\partial_{ab}u
             +q\,\mathcal Q(g_\eta,Du,Du)
             +Dq\cdot\mathcal B(g_\eta)Du,              \tag{1.3}
\] with smooth bounded algebraic coefficients while \(|u|\) is small.
Set \[
 L=\partial_t-qg_\eta^{ab}\partial_{ab},\qquad
 v=|u|_e^2,\qquad w=q|Du|_e^2 .
\] The fixed cutoff gives \[
 |Dq|^2/q+|D^2q|\le C_\chi
\] uniformly in \(\eta\). Differentiating (1.3), using ellipticity near
\(e\), and absorbing mixed Hessian terms by Young's inequality gives \[
 Lv\le-cw+Cv,\qquad
 Lw\le-cq^2|D^2u|^2+Cw^2+Cw.                         \tag{1.4}
\] For clarity, differentiation of the principal coefficient contributes
\(qDq\,Du\,D^2u\) and \(q^2(Du)^2D^2u\). Their costs are bounded by a
fraction of \(q^2|D^2u|^2\) plus \(Cw+Cw^2\). The differentiated cutoff
drift contributes \(D^2q\,w\) and \((|Dq|/\sqrt q)w^{3/2}\), bounded by
the same right side. No third derivative of \(q\) or input derivative
bound is used. In the first inequality the nonlinear term is
\(O(|u|w)\); choosing \(a_0\) small absorbs it.

Fix \(T>0\) independent of the input, and take \(K>2e^{CT}\). Then
\(v\le a^2e^{Ct}\) while the solution exists. At a spatial maximum of
\(F=w/(Ka^2-v)\), the quotient gradient terms cancel, and (1.4) implies
\[
 LF\le\{C(Ka^2-v)-c\}F^2+CF
                   \le-c'F^2+CF.                    \tag{1.5}
\] After reducing \(a_0\), scalar comparison yields \[
 \|u(t)\|_\infty\le Ca,\qquad
 q|Du(t)|^2\le Ca^2/t\quad(0<t\le T).                \tag{1.6}
\] The case \(a=0\) is the stationary flat solution.

For fixed \(\eta\), the second estimate is an ordinary gradient bound on
every \([\delta,T]\), because \(q\ge\eta^2\). Together with uniform
ellipticity and small \(\|u\|_\infty\), smooth parabolic continuation
extends this solution to the same time \(T\). The constants for higher
global derivatives may depend on \(\eta\); that is harmless. On compact
subbands where \(\chi=1\), interior parabolic estimates applied to (1.3)
instead give \[
 \|g_\eta(T)-e\|_{C^2(\{15/8<|x|<17/8\})}\le Ca,     \tag{1.7}
\] uniformly in \(\eta\). One first uses (1.6) on \([T/2,T]\); rescaling
\(u\) by \(a\) gives bounded coefficients and then the usual interior
bootstrap. This explains the linear dependence on \(a\) in (1.7), not
just a bound independent of the input.

\hypertarget{keeping-the-actual-outer-collar}{%
\subsubsection{1.2 Keeping the actual outer
collar}\label{keeping-the-actual-outer-collar}}

For a fixed smooth \(g_0\), let \(H(s)\) be its ordinary smooth
Ricci--DeTurck flow on its initial existence interval. For \(\eta^2T\)
in that interval, \(h_\eta(t)=H(\eta^2t)\) tends to \(g_0\) in
\(C^\infty\), uniformly for \(0\le t\le T\). On any collar \(U\) where
\(\chi=0\), both \(g_\eta\) and \(h_\eta\) solve the same
\(\eta^2\)-speed DeTurck equation.

Their difference \(z\), initially zero, satisfies \[
 \partial_t z=\eta^2g_\eta^{ab}\partial_{ab}z
       +\eta^2\mathcal Q(g_\eta,Dz,Dz)
       +\eta^2B_{g_0}Dz+\eta^2C_{g_0}z.               \tag{1.8}
\] The last two coefficients, also depending algebraically on
\(g_\eta,z\), are bounded by constants of this fixed smooth datum,
uniformly for small \(\eta\); this follows by expanding
\(Dg_\eta=Dh_\eta+Dz\). Smallness of \(z\) absorbs its quadratic
gradient term in the equation for \(|z|^2\). Therefore \[
 (\partial_t-\eta^2g_\eta^{ab}\partial_{ab})|z|^2
                                      \le C_{g_0}\eta^2|z|^2 .
\] Exponential barriers from the two boundary components of \(U\) give,
on every smaller collar at distance \(d>0\), \[
 |z|\le C_{g_0}\exp[-c d^2/(\eta^2t)],\qquad 0<t\le T. \tag{1.9}
\] For example use barriers
\(e^{-\lambda r+C\eta^2\lambda^2t+ C_{g_0}\eta^2t}\) in the radial
distance, and optimize \(\lambda\). Terms from the radial Hessian are
absorbed by increasing \(C\).

Interior parabolic estimates on spatial scale \(\eta\sqrt t\), followed
by differentiation on smaller collars, multiply (1.9) only by powers of
\((\eta\sqrt t)^{-1}\). The rescaled coefficient bounds follow from
(1.6) and interior bootstrapping; the coefficients involving \(h_\eta\)
are smooth for this fixed datum. Thus, for each fixed integer \(m\), \[
 g_\eta\longrightarrow g_0\quad\hbox{in }C^m
 \text{ on smaller }\{\chi=0\}\text{ collars, uniformly on }[0,T].
                                                               \tag{1.10}
\] Exponential decay also handles \(t\downarrow0\). No smooth limit
through the degeneracy interface is asserted. Only collars a positive
distance from \(\operatorname{supp}\chi\) are used.

There is an independent classical check of this collar step. On a ball
inside such a buffer put \(s=\eta^2t\). The rescaled flow is the
ordinary small-\(C^0\) DeTurck flow.
\href{https://arxiv.org/pdf/0706.0421}{Schnürer--Schulze--Simon, Theorem
4.4} gives interior derivatives up to the initial time, depending on the
fixed smooth initial metric, not on boundary derivatives. Integrating
the local equation gives \(O_{g_0}(\eta^2T)\) convergence in each fixed
interior \(C^m\) norm. This checks the uniform-in-time version of
(1.10); it is not used to assert a global smooth degenerate flow.

\hypertarget{critical-energy-is-retained-in-the-annulus}{%
\subsubsection{1.3 Critical energy is retained in the
annulus}\label{critical-energy-is-retained-in-the-annulus}}

Put \(E_\eta(t)=\int_\Omega|\operatorname{Rm}(g_\eta)|^2\). The
geometric curvature variation for \(-2q\operatorname{Ric}\) has the
structure \[
 q\Delta\operatorname{Rm}
 +\nabla q*\nabla\operatorname{Rm}
 +\nabla^2q*\operatorname{Rm}
 +q\operatorname{Rm}*\operatorname{Rm}.
\] It is the localized Ricci curvature equation of
\href{https://www.numdam.org/item/10.24033/asens.1644.pdf}{Yang, Section
9, equation (9.7)}. Here only that equation is used, not his theorem's
additional \(L^p\) Ricci hypothesis. The Lie derivative in (1.2)
contributes only boundary transport to the integral.

Integrate the \(\nabla^2q\) term by parts. Metric compatibility then
gives, with \(D_\eta=\int_\Omega q|\nabla\operatorname{Rm}|^2\), \[
 E_\eta'\le-cD_\eta+CE_\eta+
          C\int_\Omega q|\operatorname{Rm}|^3+B_\eta. \tag{1.11}
\] Every mixed derivative is absorbed using \(|\nabla q|^2/q\le C\),
which follows from the coordinate cutoff bound and zeroth-order metric
equivalence. In particular no bound on the metric Hessian of \(\chi\) is
assumed. Near \(\partial\Omega\), \(q=\eta^2\) and \(Dq=0\). Equation
(1.10), with \(m=3\), bounds all diffusion and transport fluxes there,
so \[
                         |B_\eta|\le C_{g_0}\eta^2.   \tag{1.12}
\] Large curvature of the arbitrary extension outside \(\Omega\) does
not enter this inequality.

One may alternatively estimate the Hessian term directly:
\(|\nabla_g^2q|\le C(1+a/\sqrt t)\) by (1.6). This replaces \(CE_\eta\)
in (1.11) by \(C(1+a/\sqrt t)E_\eta\). Its uniformly integrable time
coefficient gives the same conclusion below, without that integration by
parts.

The inhomogeneous Sobolev inequality on the fixed annulus is uniform
because \(g_\eta\) stays uniformly close to \(e\). Apply it and Kato's
inequality to \(\sqrt q|\operatorname{Rm}|\): \[
 \|\sqrt q|\operatorname{Rm}|\|_4^2
        \le C(D_\eta+E_\eta),\qquad
 \int_\Omega q|\operatorname{Rm}|^3
        \le C\sqrt{E_\eta}(D_\eta+E_\eta).             \tag{1.13}
\] No zero boundary values are needed in this inhomogeneous inequality.
Choose \(\varepsilon_0\) small and then \(\eta\) small for the fixed
datum. A continuity bootstrap absorbs the first term on the right of
(1.13) into (1.11). Gronwall gives \[
                     E_\eta(T)\le CE_0+C_{g_0}\eta^2. \tag{1.14}
\] Both the time and \(C\) are independent of input derivatives.

\hypertarget{finishing-the-smooth-cap}{%
\subsubsection{1.4 Finishing the smooth
cap}\label{finishing-the-smooth-cap}}

On \(15/8<|x|<17/8\), interpolate from \(e\) on the inside to
\(g_\eta(T)\) on the outside using a fixed smooth cutoff. Equation (1.7)
bounds its full curvature cost by \(Ca^2\). Use \(e\) further inside and
\(g_\eta(T)\) up to radius \(27/8\).

On \(27/8<|x|<7/2\), interpolate back to the original \(g\), then use
\(g\) for \(7/2<|x|<4\). This second collar lies in the region of
(1.10). Its metric converges in \(C^2\) to the original one for this
fixed datum, so its curvature integral equals the original collar
integral plus \(o_\eta(1)\). Convex interpolation preserves positivity.

The resulting smooth cap has total cost \[
 C(E_0+a^2)+C_{g_0}\eta^2+o_\eta(1).
\] If \(E_0+a^2>0\), choose \(\eta\) so that the last errors are at most
\(E_0+a^2\). This choice may depend on every derivative of the
individual input, but the final bound does not. If both quantities are
zero, the input is \(e\) and the flat cap suffices. This proves Theorem
1.

\hypertarget{near-einstein-energy-supplies-the-global-curvature-budget}{%
\subsection{2. Near-Einstein energy supplies the global curvature
budget}\label{near-einstein-energy-supplies-the-global-curvature-budget}}

For a closed oriented four-manifold set \[
 \mathcal E(g)=\inf_{\lambda\ne0}
   \frac{\int|\operatorname{Ric}_g-\lambda g|^2\,dV_g}{\lambda^2V_g},
 \qquad V_g=\operatorname{Vol}_g(M),
\] and write \[
 P=\int|\operatorname{Ric}^{\circ}|^2,\quad
 S=\int s^2,\quad R=P+S/4,\quad
 L_s=S-\frac{(\int s)^2}{V}\ge0 .
\] In the actual positive-budget family \(R>0\), and
\(\mathcal E=(4P+L_s)/R\). Hence \(P\le\mathcal E R/4\). The
Chern--Gauss--Bonnet identity gives \[
 8\pi^2\chi
 =\int(|W^+|^2+|W^-|^2)+R/6-2P/3
 \ge(1-\mathcal E)R/6.                               \tag{2.1}
\] Consequently, if \(\mathcal E<1/2\), \[
 R\le96\pi^2\chi,\qquad
                   \int|\operatorname{Rm}|^2\le C\chi. \tag{2.2}
\] The second estimate follows by inserting this bound into the same
identity and the curvature decomposition. These bounds need neither
scalar normalization nor a volume estimate.

On the balanced actual family \[
 M_j=X_j\# k_j\overline{\mathbb {CP}}^{\,2}
           \#\ell_j(S^1\times S^3),\quad
 A_j=c_1^2(X_j),\quad Q_j=k_j+4\ell_j,\quad 0<h_j=Q_j/A_j\to0,
\] one has \(\chi(M_j)=A_j/3+k_j-2\ell_j=O(A_j)\). Thus every sequence
with bounded \(\mathcal E/h_j\) automatically has full curvature energy
\(O(A_j)\).

Here \(X_j\) has its complex orientation and is a finite cover of the
fixed closed complex-hyperbolic surface used in I44/I52. The blowup and
handle class supply is the inherited actual one of I03, not a formal
vector on a cap. Its two classical all-metric budgets used below are,
explicitly, \[
             S\ge32\pi^2A,\qquad R\ge8\pi^2A(1+h).
\] They are the scalar and Ricci-energy inequalities detailed in
\href{20260927_covering_geometry_and_sharp_subfamilies.md}{I44, Sections
5--6}, with classical input from
\href{https://arxiv.org/html/math/0003068}{LeBrun}. No monopole supply
is demanded on the auxiliary capped manifold.

\hypertarget{long-marked-annuli-force-the-optimal-first-order-lower-bound}{%
\subsection{3. Long marked annuli force the optimal first-order lower
bound}\label{long-marked-annuli-force-the-optimal-first-order-lower-bound}}

\textbf{Theorem 2 (annular first-order stability).} On the actual family
above, suppose that around every marked blowup core there are \(L_j\)
disjoint fixed-ratio annular bands, all core families mutually disjoint.
Each band has marked coordinates, identifying its filled inside with the
corresponding blowdown ball. It encloses only its designated blowup, not
a handle or another core. Assume, after scaling each band to
\(\mathcal A\), that \[
 \max_{i,\nu}\|g_j-e\|_{C^0(\mathcal A_{i,\nu,j})}
          =a_j\longrightarrow0,\qquad
                         L_jh_j\longrightarrow\infty. \tag{3.1}
\] Then \[
                       \liminf_j\frac{\mathcal E(g_j)}{h_j}\ge4.
                                                               \tag{3.2}
\] There is no scalar, gauge, derivative, or stationary assumption. Away
from these annuli the metric is arbitrary.

\textbf{Proof.} If the ratio has no bounded subsequence there is nothing
to prove. On any bounded subsequence, (2.2) supplies total curvature
energy \(O(A_j)\). Choose a minimum-energy band for each blowup. Their
total energy satisfies \[
                         \sum_i E_{i,j}\le CA_j/L_j.   \tag{3.3}
\] Call a selected band bad if \(E_{i,j}>\varepsilon_0\), and let
\(b_j\) be the number of bad blowups. Then
\(b_j\le CA_j/(L_j\varepsilon_0)=o(Q_j)\). Theorem 1 applies to every
good band. Replace only those insides by matched balls; all handles and
bad blowups remain unchanged. The resulting closed manifold is
\(N_j=X_j\#b_j\overline{\mathbb{CP}}^{\,2}  \#\ell_j(S^1\times S^3)\).
The sum of self-dual Weyl cap costs is \[
 C_{+,j}\le C\left(A_j/L_j+k_j a_j^2\right)=o(Q_j).    \tag{3.4}
\] Indeed \(A_j/(L_jQ_j)=1/(L_jh_j)\to0\) and \(k_j\le Q_j\). No maximum
individual energy needs to tend to zero.

The partial-blowdown version of
\href{20260927_covering_geometry_and_sharp_subfamilies.md}{I44, Theorem
5} retains the cost of the \(b_j\) unremoved blowups. Suppressing \(j\),
the signature formula and the original scalar identity give \[
 \begin{aligned}
 B_M&=12\pi^2(k-b)+B_N+K_+-C_+,\\
 P_M&=8\pi^2(Q-6b)+(S_M-S_0)/12+4(B_N+K_+-C_+),
 \end{aligned}                                             \tag{3.5}
\] where \(S_0=32\pi^2A\), \(B_N=\int_N|W_N^-|^2\), and \(K_+\) is the
positive-Weyl energy of the removed cores. Consequently \[
\begin{aligned}
 &(1+h)R_M\left(\mathcal E_M-\frac{4h}{1+h}\right)\\
 &\quad=\frac{1-3h}{3}(S_M-S_0)+16(B_N+K_+-C_+)
                 +(1+h)L_s-192\pi^2b .
\end{aligned}                                             \tag{3.6}
\] The actual scalar/projection budget on \(M\) gives \(S_M\ge S_0\) and
\(R_M\ge8\pi^2A(1+h)\). For \(h\le1/3\) all terms except the two
displayed costs are nonnegative. Hence \[
 \frac{\mathcal E(g_j)}{h_j}\ge\frac4{1+h_j}
       -\frac{24b_j+2C_{+,j}/\pi^2}{Q_j(1+h_j)^2}.     \tag{3.7}
\] The two costs are \(o(Q_j)\), proving (3.2). No conformal reduction
occurred. The cap is an auxiliary smooth metric, so it need not carry a
new monopole solution. The original actual supplier remains on \(M_j\),
exactly as required by I44. QED.

The proof gives a finite, all-metric estimate, not only a sequential
statement. For \(0<h\le1/7\), \(a\le a_0\), \(L\ge1\), and \(\chi/A\)
bounded as on this family, \[
 \boxed{\mathcal E(g)\ge\frac{4h}{1+h}
                -C\left(\frac1L+\frac{k}{A}a^2\right).} \tag{3.8}
\] If \(\mathcal E<1/2\), the proof just given applies with
\(b\le CA/(L\varepsilon_0)\), whether or not \(A/L\) is small. If
\(\mathcal E\ge1/2\), (3.8) follows immediately from \(4h/(1+h)\le1/2\).
This explains why the competitive-energy premise is not an assumption of
the result.

There is a useful robustness form. Leave \(r\) marked blowups without
any prescribed annular family, and let \(a_i\) be the maximal amplitude
over the \(L\) bands of each remaining blowup. Discard also those with
\(a_i>a_0\), at most \(a_0^{-2}\sum_i a_i^2\) of them. The same
partial-cap proof gives \[
 \mathcal E(g)\ge\frac{4h}{1+h}
       -C\left(\frac rA+\frac1L+\frac1A\sum_i a_i^2\right).
                                                               \tag{3.8a}
\] Thus \(r=o(Q)\), \(\sum_i a_i^2=o(Q)\), and \(Lh\to\infty\) suffice
for (3.2). A maximum amplitude tending to zero at every core is not
necessary. The unstructured regions are retained, not silently filled,
and their number is charged in the identity.

One can average amplitudes over the bands as well. If \(a_{i\nu}\) is
the amplitude of band \(\nu\) at core \(i\), select the band minimizing
\(E_{i\nu}+a_{i\nu}^2\). The sum of selected scores is at most
\(CA/L+L^{-1}\sum_{i,\nu}a_{i\nu}^2\). Discard scores larger than
\(\min\{\varepsilon_0,a_0^2\}\); the others satisfy both cap hypotheses.
Thus (3.8a) improves to \[
 \mathcal E(g)\ge\frac{4h}{1+h}
       -C\left(\frac rA+\frac1L+
                      \frac1{AL}\sum_{i,\nu}a_{i\nu}^2\right).
                                                               \tag{3.8b}
\] This is the same partial-cap argument, not a regularity assertion for
the discarded bands.

\textbf{Corollary 3 (arbitrary \(C^0\)-small perturbations of selected
recoveries).} The CH shrinking-surgery bases can be selected so that
every sequence of smooth metrics \(g_j\) satisfying \[
                   \|g_j-g_j^0\|_{L^\infty(g_j^0)}\to0
\] obeys (3.2). This removes the TT, Bianchi, constant-scalar and
amplitude-rate assumptions from the lower-bound part of I52.

\textbf{Proof.} At each fixed \(j\), choose \(L_j\) with
\(L_jh_j\to\infty\). Inside each disjoint CH blowup chart choose a
finite annular range of that many bands, with outer radius so small that
the scaled CH deviation on every band is at most a prescribed
\(\varepsilon_j\to0\). Keep its innermost radius positive at this fixed
\(j\), and then take the actual donor scale sufficiently smaller. The
exterior Yamabe convergence of I52 is uniform on this fixed finite
union, so the normalized recovery is Euclidean-close there. Choose the
donor scales to make that additional error at most \(\varepsilon_j\),
along with all previous recovery selections. Now take the outer
\(j\)-diagonal.

This order of choices matters. No control of the Yamabe factor on every
rescaled donor is inferred from exterior convergence. For the stated
competitors, the annular amplitude in (3.1) is at most
\(C(\varepsilon_j+\|g_j-g_j^0\|_\infty)\), which tends to zero. Theorem
2 applies. QED.

The statement is invariant under simultaneous pullback. Thus a
competitor which can be made \(C^0\)-close to these bases by any
orientation-preserving diffeomorphism is also covered; no bound on that
diffeomorphism is needed. The compact moving-center selection from the
companion proof remains available by choosing the finite annular ranges
uniformly at each fixed compact configuration set. The old exact-slice
quadratic remainder and equality characterization remain separate
stronger conclusions inside their specified slices.

\hypertarget{the-minimum-over-every-shrinking-zeroth-order-neighborhood}{%
\subsubsection{3.1 The minimum over every shrinking zeroth-order
neighborhood}\label{the-minimum-over-every-shrinking-zeroth-order-neighborhood}}

This extension determines an actual variational infimum, not only the
energy of a prescribed perturbation. Let \(b_j\to0\) be the maximum
Euclidean error of the selected base on its marked annuli. For small
fixed \(\rho_0>0\), define \[
 m_j(\rho)=\inf\left\{\frac{\mathcal E(g)}{h_j}:
       g\text{ smooth},\
       \|g-g_j^0\|_{L^\infty(g_j^0)}\le\rho\right\},
 \qquad 0\le\rho\le\rho_0 .
\] The amplitude on each band is at most \(b_j+(1+b_j)\rho\). Equation
(3.8) and \(k_j\le Q_j\) give \[
 \frac4{1+h_j}
   -C\left((L_jh_j)^{-1}+(b_j+\rho)^2\right)
 \le m_j(\rho)\le\frac{\mathcal E(g_j^0)}{h_j}.        \tag{3.9}
\] The right side tends to four by the actual recovery construction. It
follows for every radius sequence \(\rho_j\to0\), with no rate relative
to \(h_j\), that \[
                         \lim_j m_j(\rho_j)=4.        \tag{3.10}
\] Also both iterated lower/upper limits as \(\rho\downarrow0\) equal
four. No existence of the inner \(j\)-limit for a fixed positive
\(\rho\) is asserted: (3.9) gives only the interval \([4-C\rho^2,4]\) in
that case.

For the intrinsic distance \[
 d_j(g)=\inf_{\varphi\in\operatorname{Diff}^+(M_j)}
       \|\varphi^*g-g_j^0\|_{L^\infty(g_j^0)},
\] approximate the infimum by diffeomorphisms and use invariance of
\(\mathcal E\). The same lower bound holds with \(\rho=d_j(g)\) whenever
this distance is small. In particular, for every fixed \(\delta>0\), a
subsequence satisfying \(\mathcal E(g_j)/h_j\le4-\delta\) must
eventually satisfy \[
                    d_j(g_j)\ge c\min\{\rho_0,\sqrt\delta\}>0.
                                                               \tag{3.11}
\] The vanishing \(b_j,(L_jh_j)^{-1},h_j\) errors are absorbed first.
This is separation from the selected recovery orbits, not a
classification of such a putative new geometry. It also does not prove
that every coefficient-four sequence approaches those orbits. The
quadratic error in (3.9) is a proved bound, not a claim of optimal
amplitude dependence.

\hypertarget{what-changed-and-what-remains}{%
\subsection{4. What changed and what
remains}\label{what-changed-and-what-remains}}

The matched cap is stronger than the previously requested vanishing
modulus: its cost is quadratic in the zeroth-order amplitude plus the
actual curvature energy. The localized flow does not need to preserve
the gauge fields or converge to an Einstein metric. It constructs a
temporary filling metric used in an exact topological identity. This
distinction avoids imposing unnecessary flow conditions.

The fixed-base exact-gauge obstruction in X09 therefore does not block
the first-order estimate. The intrinsic cap argument covers arbitrary
\(C^0\)-nearby competitors without solving that gauge. The unrestricted
infimum over every metric on \(M_j\) is still not computed: no theorem
here supplies the long Euclidean-close annuli for all near-minimizers.
The general weak zero-Weyl endpoint is also unchanged.

Classical inputs are DeTurck's strictly parabolic formulation, smooth
parabolic existence and interior estimates, Yang's localized curvature
evolution, and the Sobolev inequality on a fixed annulus. The localized
continuation, collar freezing, matching construction, and their
application to the actual surgery problem are proved above. This record
makes no worldwide priority claim.

\end{document}
